% =============================================================================
% AGENT   : Statistician + Geneticist
% MODULE  : Material and Methods — English
% =============================================================================

\section{Material and Methods}
\label{sec:methods}

\subsection{Marker Data Simulation}
\label{sec:sim}

To demonstrate the analytical pipeline, a dataset of 60 molecular markers
was simulated using the Python programming language (v3.10+) with the
\texttt{pandas} (v2.x) and \texttt{numpy} (v1.24+) libraries.
The simulation comprised 36 \snp{} markers and 24 \ssr{} markers.

Major allele frequencies ($p$) for \snp{} markers were sampled from a
Beta distribution with shape parameters $\alpha = 6$, $\beta = 2$, yielding
allele frequencies skewed towards common variants, consistent with
empirical observations in crop genomes \autocite{Elshire2011}. For
\ssr{} markers, frequencies were drawn from a symmetric Beta distribution
($\alpha = \beta = 3$), reflecting the higher polymorphism expected for
multiallelic loci modelled as biallelic equivalents.

\subsection{Statistical Parameters}
\label{sec:stats}

For each marker, the following statistics were computed:

\paragraph{Polymorphism Information Content.}
Following \textcite{Botstein1980}, \pic{} was calculated as:
\begin{equation}
    \text{PIC} = 1 - \left(p^2 + q^2\right) - 2p^2q^2
    \label{eq:pic}
\end{equation}
where $p$ and $q = 1 - p$ are the frequencies of the two alleles.

\paragraph{Expected Heterozygosity.}
Under Hardy--Weinberg equilibrium:
\begin{equation}
    H_e = 2pq
    \label{eq:he}
\end{equation}

\paragraph{Observed Heterozygosity.}
Simulated by adding Gaussian noise $\varepsilon \sim \mathcal{N}(0, 0.04)$
to $H_e$, clipped to $[0, 1]$:
\begin{equation}
    H_o = \text{clip}\left(H_e + \varepsilon,\; 0,\; 1\right)
    \label{eq:ho}
\end{equation}

\subsection{Data Visualisation}
\label{sec:viz}

Boxplots comparing \pic{} and \he{} between \snp{} and \ssr{} marker
classes were generated with the \texttt{matplotlib} library (v3.7+) and
exported as PDF files for direct inclusion in this document. All analyses
are fully reproducible via the script \texttt{scripts/plot\_test.py}
(random seed: 42).
